\documentclass[runningheads]{llncs}

\usepackage{amsmath,amssymb,amsfonts}
\usepackage{algorithm}
\usepackage{algpseudocode}
\usepackage{booktabs}
\usepackage{multirow}
\usepackage{graphicx}
\usepackage{xcolor}
\usepackage{url}
\emergencystretch=2em

\newcommand{\method}{FedProtoAug}
\newcommand{\clip}{CLIP}
\newcommand{\R}{\mathbb{R}}
\newcommand{\norm}[1]{\left\lVert #1 \right\rVert}

\begin{document}

\title{\method: Lightweight Prototype-Guided Feature Augmentation for Communication-Efficient Federated Learning under Extreme Label Skew}

\author{Anonymous Authors}
\institute{Anonymous Institution}

\maketitle

\begin{abstract}
Federated learning with pretrained vision-language encoders offers a practical route to strong visual recognition without exchanging raw data or fine-tuning large backbones. However, under extreme label skew, many clients observe only a small subset of classes, so the local classifier receives little or no supervision for absent categories. Recent feature-generation approaches such as GGEUR model full class-wise covariance geometry, but this requires transmitting high-dimensional second-order statistics. We present \method, a lightweight prototype-guided feature augmentation method for frozen \clip{} feature spaces. Each client uploads only class prototypes and counts; the server aggregates global prototypes; clients synthesize features on demand for weak or missing classes by sampling near the global prototype and optionally refining samples with a conditional residual MLP. On CIFAR-100 with Dirichlet $\alpha=0.01$, \method{} achieves $77.26\%\pm0.21$ accuracy, a gain of over 12 points above the covariance-based GGEUR baseline recorded in our baseline workbook; on Tiny-ImageNet with 200 classes, it reaches $72.65\%\pm0.08$, exceeding the recorded GGEUR baseline by more than 13 points. Across severe label-skew settings, \method{} improves accuracy over the provided covariance-based baselines while reducing statistic communication by about $500\times$.
\end{abstract}

\keywords{Federated learning \and Label skew \and Feature augmentation \and Prototype aggregation \and Pretrained models \and Communication efficiency}

\section{Introduction}

Federated learning (FL) trains a model over decentralized clients while keeping raw data local \cite{mcmahan2017fedavg}. A central challenge is statistical heterogeneity: client data distributions often differ in both labels and feature appearance. The challenge becomes severe under label skew, where each client observes only a few classes. In this regime, local training produces classifiers biased toward locally present labels, and standard model averaging cannot recover class boundaries that many clients never update.

Large pretrained encoders change the structure of this problem. When a frozen \clip{} encoder \cite{radford2021clip} maps images into a semantically organized feature space, the main bottleneck under extreme Dirichlet label skew is not necessarily fine-grained feature shift; it is class absence. If a client lacks a class, the local linear head has no positive examples for that class, even though a global class prototype can be estimated from other clients. This observation motivates a simple question: how much of the benefit of heavy distribution-alignment methods can be recovered by transmitting only class means?

This question is important because frozen feature FL is a common deployment pattern. A hospital, mobile application, or edge fleet may already rely on a foundation encoder that is too expensive or too sensitive to fine-tune on every client. In that setting, the trainable part is often only a linear or shallow classifier. Classical non-IID FL methods still help, but they do not directly solve the fact that absent labels produce absent gradients. The classifier's row for a missing class is updated only through negative examples, so the class can be pushed away even when the global model should preserve it.

Recent work GGEUR \cite{ma2025ggeur} addresses heterogeneity by estimating global geometric shapes of embedding distributions and using them to guide local generation. This is effective, but in frozen 512-dimensional feature spaces it requires transmitting class-wise covariance matrices, whose cost scales as $O(Cd^2)$ per client. For CIFAR-100 with $d=512$, this is roughly $100$ MB per client per round for the geometric statistics alone. By contrast, class prototypes scale as $O(Cd)$ and cost roughly $0.2$ MB per client per round.

We propose \method, a communication-efficient method for extreme label skew with frozen pretrained features. Clients compute local class prototypes and counts. The server aggregates them into global class prototypes. During local training, each client detects weak classes and augments them by sampling normalized features around the corresponding global prototype. A lightweight conditional residual MLP can further refine the sampled features, but the core augmentation requires no covariance matrices, no raw data transfer, and no backbone updates.

The method is deliberately conservative. It does not attempt to reconstruct the full global distribution, and it does not assume that the server can synthesize raw images. Instead, it exploits a property of strong normalized feature spaces: class means are often stable semantic anchors. A synthetic feature drawn near such an anchor is not intended to be a perfect sample from the real class distribution. It only needs to provide a useful positive update to the local classifier for a class that would otherwise be absent. This distinction makes the approach lightweight and explains why it is particularly effective in extreme label skew.

Our contributions are:
\begin{itemize}
    \item We identify class absence as a primary bottleneck in frozen pretrained FL under extreme label skew and show that prototype-guided feature augmentation is sufficient to recover most of the lost accuracy.
    \item We introduce \method, a lightweight client-side feature augmentation framework based on prototype aggregation, demand-driven weak-class synthesis, and an optional conditional MLP refiner.
    \item We provide a communication analysis showing about $500\times$ lower statistic communication than full covariance-based GGEUR under CIFAR-100-scale settings.
    \item We report experiments on CIFAR-10, CIFAR-100, and Tiny-ImageNet. FedFM/FedProtoAug results use the provided experiment document, baseline results use the provided GGEUR workbook when available, and missing entries are supplemented from GGEUR Tables 1--3.
\end{itemize}

\section{Related Work}

\subsection{Federated Learning under Heterogeneity}
FedAvg \cite{mcmahan2017fedavg} averages local model parameters and remains the standard FL baseline. Many methods improve stability under non-IID data. FedProx \cite{li2020fedprox} adds a proximal term to restrict local drift. SCAFFOLD \cite{karimireddy2020scaffold} uses control variates to reduce client drift. FedDyn \cite{acar2021feddyn} introduces dynamic regularization, and MOON \cite{li2021moon} uses model-contrastive learning. These methods mainly modify optimization or aggregation, whereas \method{} focuses on the missing-supervision problem in the frozen feature space.

\subsection{Federated Data and Feature Augmentation}
Data augmentation has been explored as a direct response to client bias. FedMix \cite{yoon2021fedmix} extends Mixup through mean-augmented federated learning. FedGen \cite{zhu2021fedgen} trains a server-side generator for data-free distillation. FedFA \cite{zhou2023fedfa} manipulates latent feature statistics to regularize feature shift during model training. GGEUR \cite{ma2025ggeur} uses global geometric shapes to locally simulate global embedding distributions and obtains strong performance under label and domain skew. \method{} differs in its minimalist statistic design: it uses class prototypes and counts rather than full covariance geometry.

\subsection{Prototype-Based Federated Learning}
Prototype communication is an established way to exchange compact semantic information. FedProto \cite{tan2022fedproto} aggregates class prototypes and regularizes local representations toward global prototypes. Recent generative prototype approaches such as GFPL \cite{lu2026gfpl} further use prototype distributions to generate pseudo-features under imbalance. \method{} is closest in spirit to prototype communication but targets frozen pretrained feature extractors and uses prototypes directly for demand-driven local augmentation rather than representation regularization.

GFPL and \method{} differ in their assumptions and target regimes. GFPL addresses resource-constrained and data-imbalanced federated vision tasks with a generative prototype-learning framework, whereas \method{} focuses on frozen pretrained feature spaces and extreme label skew. In \method, the prototype is not a general-purpose generator input; it is a compact class anchor used to restore missing positive supervision for a linear head.

\subsection{Pretrained Vision-Language Models in FL}
Pretrained vision-language models have been adapted to FL through lightweight components. FedCLIP \cite{lu2023fedclip} studies generalization and personalization of \clip{} with adapters. FedTPG \cite{qiu2024fedtpg} learns a federated text-driven prompt generator for vision-language models. \method{} follows the same practical principle: keep the large pretrained backbone frozen and communicate only lightweight task-specific information.

\subsection{Position Relative to GGEUR}
The closest empirical reference point for this work is GGEUR \cite{ma2025ggeur}. Both methods operate in an embedding space and generate local training samples to approximate missing global information. The difference is the statistic used to define that information. GGEUR estimates class-wise geometric shape, which captures richer within-class variation but requires second-order communication. \method{} keeps only first-order class location and uses it to solve the missing-label supervision problem. Table~\ref{tab:method_compare} summarizes the design contrast.

\begin{table}[t]
\centering
\caption{Method-level comparison with GGEUR in a frozen 512-dimensional feature space.}
\label{tab:method_compare}
\begin{tabular}{lcc}
\toprule
Aspect & \method{} & GGEUR \\
\midrule
Statistic & Mean prototype & Covariance geometry \\
Scaling & $O(Cd)$ & $O(Cd^2)$ \\
Client synthesis & Prototype Gaussian & Geometric directions \\
Backbone & Frozen & Frozen VFM setting \\
Target regime & Class absence & Distribution mismatch \\
Main trade-off & Low cost & Richer shape \\
\bottomrule
\end{tabular}
\end{table}

This comparison also motivates our experimental framing. We do not claim that prototypes dominate covariance information in every regime. Instead, we argue that under extreme label skew and a strong frozen encoder, class absence is a more urgent failure mode than detailed within-class geometry. When the regime changes, the preferred point on the communication-accuracy curve can also change.

\section{Method}

\subsection{Problem Formulation}
We consider $K$ clients and $C$ classes. Client $k$ owns a private dataset $\mathcal{D}_k=\{(x_i,y_i)\}$ sampled from a non-IID distribution. A frozen pretrained encoder $f:\mathcal{X}\rightarrow\R^d$ maps each image to a $d$-dimensional feature, with $d=512$ for the \clip{} ViT-B/16 setting used in our experiments. Only a classifier head $h_\theta:\R^d\rightarrow\R^C$ is trained federatively.

The goal is to learn a global head $\theta$ without sharing raw images or per-sample features. Under extreme Dirichlet label skew, many clients have $n_k^c=0$ or very small $n_k^c$ for class $c$. \method{} addresses this by providing local synthetic features for weak classes using global prototype information.

\subsection{Prototype Aggregation}
Each client computes a local class prototype for classes with enough local samples:
\begin{equation}
\mu_k^c = \frac{1}{n_k^c}\sum_{(x_i,y_i)\in\mathcal{D}_k, y_i=c} \frac{f(x_i)}{\norm{f(x_i)}_2}.
\end{equation}
The client uploads $(\mu_k^c,n_k^c)$, and the server aggregates prototypes by sample-count weighting:
\begin{equation}
\mu^c = \frac{\sum_{k=1}^{K} n_k^c \mu_k^c}{\sum_{k=1}^{K} n_k^c}.
\end{equation}
For multi-modal local class distributions, the same rule can be applied per mode after client-side clustering; the implementation supports this option, while our main paper emphasizes the single-prototype setting because it is the cleanest and most communication-efficient variant.

The minimum sample threshold prevents unreliable local means from dominating the global statistic. If a client has only one or two examples of a class, its local prototype may be noisy. In the implementation, such classes can be skipped during statistic upload while still being augmented during local training if the server has a global prototype estimated from other clients. This asymmetry is useful: a weak local class should receive help, but it should not necessarily contribute a high-variance estimate to the global prototype.

Aggregation uses only sample counts and normalized prototypes. Therefore, it is invariant to local batch ordering and independent of the local classifier state. This matters in our setting because the classifier head is updated every round, while frozen features and class means remain stable enough to serve as a slowly changing global memory. The server can update prototypes every round or less frequently; in our experiments, prototype upload is performed with the normal communication loop.

\subsection{Demand-Driven Weak-Class Augmentation}
A client marks class $c$ as weak if its local count satisfies either
\begin{equation}
n_k^c < n_{\min}
\end{equation}
or if $n_k^c$ falls in the bottom $p$ percentile among all class counts on that client. The number of generated features can be set dynamically as
\begin{equation}
m_k^c = \min\{\lambda(n^\star-n_k^c)_+, m_{\max}\},
\end{equation}
or fixed per weak class in the implementation. Synthetic features are sampled near the global prototype:
\begin{equation}
z_0 = \operatorname{norm}(\mu^c + \epsilon), \quad \epsilon\sim\mathcal{N}(0,\sigma^2 I).
\end{equation}
The synthetic set is concatenated with the client's real frozen features, and the local classifier head is trained with cross-entropy. This process gives absent or rare classes positive supervision without transmitting examples.

Demand-driven synthesis avoids a common failure mode of indiscriminate augmentation. If a client already has many examples of class $c$, additional prototype-centered synthetic features may over-smooth the local decision boundary. We therefore generate features only for missing or underrepresented classes. This keeps the local training set close to the real client distribution while supplying the supervision needed to prevent class collapse.

The proposal variance $\sigma$ controls the diversity-fidelity trade-off. A small $\sigma$ produces features that act almost like repeated class prototypes; a large $\sigma$ may generate ambiguous points that cross class boundaries. We use $\sigma=0.05$ after feature normalization, which empirically produces enough diversity for classifier training without requiring covariance estimation. Because the head is linear, even prototype-near samples can provide meaningful class-row updates.

\subsection{Conditional MLP Refiner}
The optional refiner is a conditional residual MLP. Given a proposal $z_0$ and a class embedding $e_c$, it predicts a residual:
\begin{equation}
\delta = g_\phi([z_0; e_c]), \quad z = \operatorname{norm}(z_0+\delta).
\end{equation}
In our MLP configuration, $e_c\in\R^d$, and $g_\phi$ uses three fully connected layers with hidden width 512. The refiner is trained on the server using aggregated prototype statistics. The primary objective constrains generated features to remain close to the corresponding prototype:
\begin{equation}
\mathcal{L}_{proto} = \mathbb{E}_{c,z}\min_{\mu\in\mathcal{P}^c}\norm{z-\mu}_2^2,
\end{equation}
where $\mathcal{P}^c$ is the set of aggregated prototypes or modes for class $c$. The codebase also supports protoflow and diagonal-covariance matching losses:
\begin{equation}
\mathcal{L} = w_{protoflow}\mathcal{L}_{protoflow}+w_{proto}\mathcal{L}_{proto}+w_{cov}\mathcal{L}_{cov}+w_{reg}\mathcal{L}_{reg}.
\end{equation}
Our ablations show that protoflow matching contributes little in the frozen \clip{} feature space, so the paper's main interpretation is prototype-guided augmentation rather than histogram matching.

The residual form is important. Directly mapping noise and labels to features would require the refiner to learn the entire feature manifold. Instead, the proposal already lies near the intended semantic region, and the MLP only adjusts it. This reduces the risk that the refiner produces off-manifold artifacts. After refinement, features are normalized again to match the \clip{} feature convention and the classifier's training distribution.

In practice, we pretrain the refiner once from initial aggregated statistics and optionally fine-tune it every few communication rounds. The refiner is not a privacy-critical component in the same sense as raw data: it is trained only on aggregated prototypes and protoflowes, and it is broadcast to clients like any other model parameter. When communication is severely constrained, the refiner can be omitted entirely; Table~\ref{tab:ablation} shows that prototype sampling alone already provides most of the accuracy gain.

\subsection{Privacy and Security Scope}
\method{} does not provide a formal differential privacy guarantee. It reduces exposure by transmitting class-level aggregate statistics rather than raw images, raw per-sample features, or full feature sets. The prototype of a class with many samples is a coarse summary, but a prototype computed from very few samples could leak more information. For this reason, the method includes a minimum sample threshold for statistic upload. If a deployment requires stronger guarantees, standard secure aggregation or noise addition can be applied to the prototype sums and counts before server aggregation. These mechanisms are orthogonal to the augmentation strategy.

Compared with covariance-based communication, prototypes also expose less high-order information about local distributions. A covariance matrix reveals directions of local variation and can be sensitive when classes have few examples. \method{} intentionally avoids such second-order statistics in the main variant. This choice is motivated by both communication efficiency and the observation that class means are sufficient for the target extreme label-skew regime.
If Gaussian noise is added for differential privacy, the perturbation is applied to $O(Cd)$ prototype coordinates rather than $O(Cd^2)$ covariance entries. For the same privacy budget, this lower-dimensional statistic is easier to stabilize and communicate, giving prototype aggregation an additional practical advantage under privacy constraints.

\subsection{Optimization Intuition}
Consider a linear classifier with weights $W=[w_1,\ldots,w_C]^\top$ and a local empirical loss on client $k$. If class $c$ is absent from $\mathcal{D}_k$, there is no positive term that increases the logit $w_c^\top z$ for class-$c$ features. For any local sample from another class, the cross-entropy gradient contributes a negative-class update to $w_c$. Accumulated over several local epochs, this can move $w_c$ away from the global direction needed for class $c$.

Prototype augmentation changes this gradient structure. A synthetic pair $(z,c)$ with $z$ near $\mu^c$ contributes
\begin{equation}
\nabla_{w_c}\ell(W;z,c) = (p_c(z)-1)z,
\end{equation}
where $p_c(z)$ is the predicted probability for class $c$. Since $p_c(z)<1$ unless the classifier is already confident, the update pulls $w_c$ toward $z$ and therefore toward $\mu^c$. For absent classes, this creates the positive signal that local training lacks. For rare classes, it stabilizes the local estimate without requiring the client to own many examples.

This argument does not require $z$ to be an exact sample from the true class distribution. It only requires $z$ to lie in a region that should score highly for class $c$. Frozen \clip{} features make this condition plausible because normalized class clusters tend to be semantically organized. Thus, \method{} can be effective even with an isotropic Gaussian proposal.

\begin{algorithm}[t]
\caption{\method}
\label{alg:fedprotoaug}
\begin{algorithmic}[1]
\Require Clients $\{1,\ldots,K\}$, classes $C$, frozen encoder $f$, rounds $R$
\State Initialize global head $\theta^0$
\State Each client extracts frozen features $f(x)$ locally
\State Each client uploads local prototypes and counts $\{(\mu_k^c,n_k^c)\}$
\State Server aggregates global prototypes $\{\mu^c\}$
\State Optionally train conditional refiner $g_\phi$
\For{$r=1$ to $R$}
    \For{each client $k$ in parallel}
        \State Receive $\theta^{r-1}$, $\{\mu^c\}$, and optional $g_\phi$
        \State Identify weak classes by count threshold or percentile
        \For{each weak class $c$}
            \State Sample $z_0=\operatorname{norm}(\mu^c+\epsilon)$, $\epsilon\sim\mathcal{N}(0,\sigma^2I)$
            \State If refiner is enabled, set $z=\operatorname{norm}(z_0+g_\phi([z_0;e_c]))$; otherwise $z=z_0$
            \State Add $(z,c)$ to the local synthetic feature set
        \EndFor
        \State Train local head $\theta_k^r$ on real and synthetic features
        \State Upload $\theta_k^r$ and updated prototype statistics
    \EndFor
    \State Server aggregates $\theta^r=\sum_k \frac{|\mathcal{D}_k|}{\sum_j|\mathcal{D}_j|}\theta_k^r$
    \State Server updates global prototypes; optionally fine-tunes $g_\phi$
\EndFor
\State \Return Global head $\theta^R$
\end{algorithmic}
\end{algorithm}

\subsection{Communication Analysis}
For $d$-dimensional features and $C$ classes, \method{} transmits $C(d+1)$ scalars per client for prototype statistics. GGEUR-style full covariance geometry transmits $Cd^2$ floating-point values per client, excluding the classifier head. Thus, the statistic ratio is approximately
\begin{equation}
\frac{Cd^2}{C(d+1)} \approx d.
\end{equation}
For $d=512$, this is about $512\times$. On CIFAR-100, \method{} requires roughly $100\times(512+1)\times4\approx0.2$ MB per client, while covariance transmission requires $100\times512^2\times4\approx100$ MB per client.

The trainable linear head has $(d+1)C$ parameters, which is also about $0.2$ MB for CIFAR-100. Thus, adding prototype upload roughly doubles the upload size of head-only FedAvg in this setting, while covariance upload increases it by about three orders of magnitude. The optional MLP refiner has about $1.1$M parameters when $d=512$ and $C=100$, corresponding to about $4.4$ MB if transmitted in full. Since the refiner can be pretrained once or updated infrequently, its amortized cost over 100 rounds is small compared with per-round covariance communication.

Table~\ref{tab:comm_scaling} summarizes asymptotic scaling. The main difference is the statistic term: prototype communication grows linearly with feature dimension, while full covariance communication grows quadratically. This distinction becomes more important as pretrained encoders move to larger embedding dimensions.

\begin{table}[t]
\centering
\caption{Statistic communication scaling per client.}
\label{tab:comm_scaling}
\begin{tabular}{lcc}
\toprule
Statistic & Scalars & Scaling \\
\midrule
Class count only & $C$ & $O(C)$ \\
Prototype + count & $C(d+1)$ & $O(Cd)$ \\
Diagonal covariance & $C(2d+1)$ & $O(Cd)$ \\
Full covariance & $C(d^2+d+1)$ & $O(Cd^2)$ \\
\bottomrule
\end{tabular}
\end{table}

\section{Experiments}

\subsection{Setup}
We evaluate on CIFAR-10, CIFAR-100, and Tiny-ImageNet. Unless stated otherwise, \method{} uses $10$ clients, Dirichlet label partitioning with concentration $\alpha$, $100$ communication rounds, $10$ local epochs, Adam with learning rate $10^{-3}$, batch size $64$, and a frozen OpenCLIP ViT-B/16 feature extractor \cite{ilharco2021openclip} with a linear head. We report top-1 accuracy. \method{} results are taken from the provided FedFM/FedProtoAug experiment document. Baseline results are taken from the provided GGEUR workbook when available; entries absent from the workbook are supplemented from Tables 1--3 of GGEUR \cite{ma2025ggeur}.

The Dirichlet concentration parameter controls the severity of label skew. Smaller $\alpha$ values create sharper class imbalance across clients. At $\alpha=0.01$, clients can miss a large fraction of classes, which directly stresses the classifier-head rows for absent labels. At $\alpha\ge0.3$, clients observe broader label support, and the problem becomes closer to ordinary distribution mismatch. We therefore separate severe and mild regimes in the analysis.

\subsection{Implementation Details}
All image features are extracted once and cached locally. This makes the federated loop head-only and isolates the effect of feature-space augmentation. The head is initialized globally and trained locally using cross-entropy. After local training, the server performs sample-count weighted FedAvg over head parameters. Prototype statistics are recomputed from real local features rather than synthetic features, so generated samples do not feed back into the global prototype estimates.

For \method, we use a minimum class-statistic threshold of $10$, proposal standard deviation $\sigma=0.05$, weak-class percentile threshold as configured in the run script, and a conditional MLP refiner when the full method is enabled. The refiner hidden width is $512$, and the local classifier remains linear. This design keeps the trainable client-side classifier small and makes the comparison with CLIP+linear baselines direct.

The baseline workbook records reproduced CLIP+linear runs under comparable settings, with some baselines using 10 local epochs. Results supplemented from GGEUR are marked with $\dagger$. Small gaps should therefore be interpreted with care because local training details, seeds, or checkpoint choices may differ.

\subsection{Ablation Study}
Table~\ref{tab:ablation} shows the contribution of each component on CIFAR-100 under $\alpha=0.01$. Prototype-only calibration improves only modestly over FedAvg, while prototype sampling gives a much larger gain. The conditional refiner further improves accuracy, but most of the improvement comes from supplying weak classes with prototype-centered samples.

\begin{table}[t]
\centering
\caption{Ablation on CIFAR-100 with Dirichlet $\alpha=0.01$.}
\label{tab:ablation}
\begin{tabular}{lcc}
\toprule
Method & Description & Acc. (\%) \\
\midrule
FedAvg & No augmentation & 53.16 \\
Proto-Cal & Prototype head calibration & 56.94 \\
Proto-Aug & Gaussian prototype sampling & 74.62 \\
\method{} & Proto-Aug + MLP refiner & 77.18 \\
\bottomrule
\end{tabular}
\end{table}

The ablation also clarifies why \method{} should not be viewed primarily as a generative modeling method. The generated features are simple, local training aids. The large jump from FedAvg to Proto-Aug indicates that missing positive class updates are the main failure mode. The smaller gap from Proto-Aug to \method{} suggests that refining the samples can improve feature quality, but it is not the source of the main effect. This is desirable for deployment because the simplest variant remains strong when the refiner is disabled.

\subsection{Severe Label Skew on CIFAR-100}
Table~\ref{tab:cifar100_severe} compares severe label skew settings. \method{} consistently outperforms the GGEUR and FedProto baselines recorded in the workbook at $\alpha\le0.09$. The largest gap appears in the most skewed regime, supporting our claim that class absence is the dominant bottleneck there.

\begin{table}[t]
\centering
\caption{CIFAR-100 under severe label skew. Baselines are from the provided GGEUR workbook.}
\label{tab:cifar100_severe}
\begin{tabular}{lccccc}
\toprule
Method & $\alpha=0.01$ & $\alpha=0.03$ & $\alpha=0.05$ & $\alpha=0.07$ & $\alpha=0.09$ \\
\midrule
FedAvg & $46.37\pm1.58$ & 45.89 & 50.75 & 54.72 & 53.76 \\
FedProto & 68.88 & 64.12 & 63.11 & 63.86 & 62.14 \\
GGEUR & $64.82\pm0.36$ & 63.42 & 68.32 & 66.45 & 66.53 \\
\method{} & $77.26\pm0.21$ & $77.72\pm0.08$ & $77.96\pm0.02$ & $78.30\pm0.12$ & $78.09\pm0.26$ \\
\bottomrule
\end{tabular}
\end{table}

\subsection{Mild Label Skew on CIFAR-100}
Table~\ref{tab:cifar100_mild} reports milder heterogeneity. \method{} remains ahead of the recorded GGEUR baseline even at $\alpha=0.5$, although the gap relative to strong baselines such as FedFA becomes much smaller than in the severe-skew regime. This suggests that frozen \clip{} prototypes remain useful beyond the most extreme partitions, while the relative advantage should not be attributed only to label absence.

\begin{table}[t]
\centering
\caption{CIFAR-100 under mild label skew. Workbook baselines are primary; $\dagger$ denotes values supplemented from GGEUR Table 1 \cite{ma2025ggeur}.}
\label{tab:cifar100_mild}
\begin{tabular}{lccc}
\toprule
Method & $\alpha=0.1$ & $\alpha=0.3$ & $\alpha=0.5$ \\
\midrule
FedMix$^\dagger$ & 73.85 & 79.62 & 81.31 \\
FEDGEN$^\dagger$ & 73.15 & 78.97 & 81.24 \\
FedFA & 73.27 & 76.35 & 77.91 \\
FedProto & 61.10 & 58.60 & 57.56 \\
FedAvg & 57.36 & 62.22 & 73.96 \\
GGEUR & 67.63 & 73.94 & 74.44 \\
\method{} & $78.69\pm0.01$ & $78.93\pm0.05$ & $79.33\pm0.04$ \\
\bottomrule
\end{tabular}
\end{table}

The mild-regime results should be read together with the communication costs and baseline settings. One possible explanation for the persistent advantage over the recorded GGEUR values is that the workbook baselines use a different local-training configuration. Another is that frozen \clip{} features make class prototypes strong enough that first-order augmentation remains beneficial even when label coverage improves. We also note that frozen \clip{} features produce well-separated, near-linear class clusters, which may give prototype-based augmentation a structural advantage over methods that invest communication in modeling within-class spread that a linear head cannot fully exploit. The safer conclusion is therefore not that \method{} is only useful under extreme skew, but that its accuracy--communication trade-off is especially attractive there.

\subsection{Multi-Dataset Results}
Tables~\ref{tab:cifar10} and~\ref{tab:tiny} show results on CIFAR-10 and Tiny-ImageNet. CIFAR-10 is relatively saturated under frozen \clip{} features, so all strong methods perform near the mid-90\% range. Tiny-ImageNet has more classes, making class absence more pronounced; \method{} substantially exceeds the recorded GGEUR and FedProto baselines in severe settings.

\begin{table}[t]
\centering
\caption{CIFAR-10 results under label skew. Baselines are from the provided GGEUR workbook.}
\label{tab:cifar10}
\resizebox{\textwidth}{!}{%
\begin{tabular}{lcccccccc}
\toprule
Method & .01 & .03 & .05 & .07 & .09 & .1 & .3 & .5 \\
\midrule
FedAvg & $79.41\pm1.35$ & 68.95 & 86.55 & 86.85 & 89.16 & 91.00 & 93.25 & 94.67 \\
FedFA & -- & -- & -- & -- & -- & 91.93 & 93.84 & 94.24 \\
FedProto & 93.93 & 89.20 & 91.22 & 90.87 & 89.89 & 89.89 & 89.89 & 84.72 \\
GGEUR & $91.48\pm0.39$ & 90.85 & 90.23 & 91.81 & 92.11 & 94.03 & 94.45 & 95.02 \\
\method{} & $93.54\pm0.28$ & $94.09\pm0.12$ & $94.06\pm0.26$ & $94.02\pm0.58$ & $94.27\pm0.10$ & $94.66\pm0.23$ & $94.86\pm0.12$ & $95.06\pm0.27$ \\
\bottomrule
\end{tabular}}
\end{table}

The Tiny-ImageNet results are especially informative because the dataset has 200 classes. With more classes and the same number of clients, the probability of local class absence increases. Prototype-guided augmentation benefits from this setting because a single global prototype per class can be shared with all clients at linear cost. Full covariance methods become more expensive as both $C$ and $d$ grow, while the prototype method remains lightweight.

For CIFAR-10, the task has only 10 classes and \clip{} features are strong, so accuracy is already high. This reduces the room for improvement and makes results sensitive to implementation details such as seed, local epoch count, and whether best-round or final-round accuracy is reported. After correcting the FedFM records, \method{} remains competitive across the full CIFAR-10 range, including the mild skew settings $\alpha=0.3$ and $\alpha=0.5$.

\begin{table}[t]
\centering
\caption{Tiny-ImageNet results under label skew. Workbook baselines are primary; $\dagger$ denotes values supplemented from GGEUR Table 1 \cite{ma2025ggeur}.}
\label{tab:tiny}
\resizebox{\textwidth}{!}{%
\begin{tabular}{lcccccccc}
\toprule
Method & .01 & .03 & .05 & .07 & .09 & .1 & .3 & .5 \\
\midrule
FedMix$^\dagger$ & -- & -- & -- & -- & -- & 60.90 & 69.19 & 70.29 \\
FEDGEN$^\dagger$ & -- & -- & -- & -- & -- & 60.36 & 68.91 & 70.31 \\
FedFA & -- & -- & -- & -- & -- & 66.51 & 71.18 & 70.95 \\
FedProto & 63.28 & 63.41 & 62.95 & 60.02 & 60.10 & 59.17 & 60.25 & 59.52 \\
FedAvg & 40.85 & 47.34 & 44.27 & 50.17 & 54.84 & 56.39 & 64.18 & $68.61\pm0.53$ \\
GGEUR & 59.07 & 59.18 & 59.43 & 61.01 & 61.30 & 61.83 & 66.19 & $67.28\pm1.13$ \\
\method{} & $72.65\pm0.08$ & $72.84\pm0.04$ & $72.67\pm0.04$ & $72.99\pm0.27$ & $73.48\pm0.12$ & $73.18\pm0.11$ & $74.26\pm0.13$ & $74.82\pm0.06$ \\
\bottomrule
\end{tabular}}
\end{table}

\subsection{Communication Efficiency}
Table~\ref{tab:comm} compares per-round communication for CIFAR-100 with a 512-dimensional frozen feature and a linear head. FedAvg communicates only the head, but it lacks class-completion information. GGEUR communicates full covariance matrices and therefore incurs roughly two orders of magnitude more total traffic over 100 rounds. \method{} adds only prototype statistics to FedAvg, keeping total communication close to head-only FL while providing strong gains under severe label skew.

\begin{table}[t]
\centering
\caption{Approximate communication cost on CIFAR-100 with $K=10$, $d=512$, and 100 rounds.}
\label{tab:comm}
\begin{tabular}{lcc}
\toprule
Method & Upload / client / round & Total upload \\
\midrule
FedAvg & Head, $\approx0.2$ MB & $\approx0.2$ GB \\
GGEUR & Head + covariance, $\approx100$ MB & $\approx100$ GB \\
\method{} & Head + prototypes, $\approx0.4$ MB & $\approx0.4$ GB \\
\bottomrule
\end{tabular}
\end{table}

\subsection{Accuracy--Communication Trade-off}
Table~\ref{tab:tradeoff} summarizes the severe-skew trade-off at $\alpha=0.01$. The pattern is consistent across CIFAR-100 and Tiny-ImageNet: \method{} obtains a large accuracy gain over the recorded GGEUR baseline while using prototype statistics whose cost is linear in the feature dimension. On CIFAR-10, the task is close to saturation and the gap is smaller, but \method{} remains competitive without second-order communication.

\begin{table}[t]
\centering
\caption{Severe-skew trade-off at $\alpha=0.01$. Accuracy values are top-1 percentages.}
\label{tab:tradeoff}
\begin{tabular}{lccc}
\toprule
Dataset & GGEUR & \method{} & Gap \\
\midrule
CIFAR-10 & $91.48\pm0.39$ & $93.54\pm0.28$ & +2.06 \\
CIFAR-100 & $64.82\pm0.36$ & $77.26\pm0.21$ & +12.44 \\
Tiny-ImageNet & 59.07 & $72.65\pm0.08$ & +13.58 \\
\bottomrule
\end{tabular}
\end{table}

The trade-off can be interpreted in terms of what each transmitted statistic is used for. A full covariance matrix can describe anisotropic class spread and therefore supports more expressive generation. However, in the extreme-skew regime, the local classifier first needs a class anchor. Once an anchor is available, a linear head can recover a large portion of the lost accuracy. The extra information in a covariance matrix may help later, but it is not the highest-leverage statistic when many clients have no examples of a class.

This also suggests a hybrid direction. A practical system could begin with \method{} in early or highly skewed rounds, then selectively request richer statistics only for classes whose global decision boundary remains unstable. Such an adaptive protocol would keep the common case lightweight while preserving an upgrade path for difficult classes. We leave this selective statistic scheduling to future work.

\subsection{Discussion}
The experiments suggest that frozen pretrained feature spaces alter the balance between class absence and distribution mismatch. Under $\alpha=0.01$, many clients lack most classes, so a global prototype from other clients provides a stable semantic anchor for classes that are locally absent. This explains why the simple Proto-Aug variant recovers most of the improvement in Table~\ref{tab:ablation}. We also observe that protoflow or histogram matching contributes limited marginal gain in the frozen \clip{} space, indicating that class means already encode much of the information needed by a linear classifier. As $\alpha$ increases, most clients observe more classes, so within-class geometry and local distribution mismatch become more relevant. This shift explains why the relative advantage changes across severe and mild skew settings.

\subsection{Limitations}
\method{} assumes that a frozen pretrained feature space is available and that class prototypes are meaningful. If the encoder is weak, poorly aligned with the target domain, or trained on a very different modality, prototype-centered samples may not approximate useful class regions. The method also assumes a shared label space across clients. If clients have private labels or incompatible taxonomies, global class prototypes need additional alignment.

Baseline methods were evaluated under comparable but not identical seeds and local training configurations; a standardized re-evaluation under a single protocol would strengthen the comparison.

\subsection{Practical Recommendations}
The empirical pattern suggests several practical rules. First, prototype augmentation should be enabled when the estimated label coverage per client is low. A simple diagnostic is the fraction of classes with $n_k^c<n_{\min}$ on each client. If many clients have weak classes, \method{} directly addresses the missing-gradient issue. If most clients already cover most classes, prototype augmentation should be used more conservatively, because the problem may be dominated by within-class distribution mismatch.

Second, the refiner should be treated as an optional quality module rather than a mandatory component. When the pretrained encoder is strong, the prototype proposal may already be close to the useful class region. In that case, disabling the refiner reduces downlink cost and implementation complexity with little accuracy loss. When the encoder is weaker or the dataset has many visually diverse classes, the refiner can improve the proposal by learning class-conditioned residual corrections.

Third, the prototype upload threshold should not be set too low. Uploading prototypes from one or two samples can add noise and may increase privacy risk. A moderate threshold, such as 10 samples per class in our runs, is a reasonable default for CIFAR-scale datasets. For smaller datasets, secure aggregation or client-side clipping can be used to stabilize prototype sums before averaging.

Finally, \method{} can be combined with existing FL optimizers. Since the method changes the local training set rather than the aggregation rule, it is compatible with FedProx, SCAFFOLD, FedDyn, or personalized heads. This is useful in deployments where optimization drift and class absence occur together. The current draft uses FedAvg to isolate the effect of feature augmentation; stronger optimizers may further improve the absolute numbers.

\subsection{Failure Modes}
Prototype-guided augmentation can fail when the global prototype is biased. This may happen if only a small subset of clients observes a class, or if those clients contain a narrow visual mode. In that case, generated features may overrepresent the observed mode and underrepresent rare modes. Multi-prototype clustering can mitigate this issue by preserving several class anchors, but it increases communication linearly in the number of modes.

Another failure mode is class confusion in the frozen feature space. If two classes have nearby prototypes, isotropic sampling can produce ambiguous synthetic features. The linear head may then receive noisy supervision. A smaller proposal variance, class-aware rejection sampling, or nearest-neighbor filtering against other class prototypes can reduce this risk. These additions preserve the lightweight nature of \method{} because they use the same prototype table.

The method also depends on stable feature normalization. If client feature extraction differs due to preprocessing, backbone checkpoint mismatch, or numerical inconsistency, aggregated prototypes can become misaligned. For this reason, all clients should share the same frozen encoder, preprocessing pipeline, and feature normalization convention. This is easier to enforce than synchronized backbone training, but it remains an important deployment assumption.

\section{Conclusion}
We introduced \method, a lightweight prototype-guided feature augmentation method for federated learning with frozen pretrained features. By communicating only class prototypes and counts, \method{} supplies weak or missing classes with synthetic local features while avoiding full covariance transmission. Experiments across CIFAR-10, CIFAR-100, and Tiny-ImageNet show that prototype sampling recovers most of the performance lost to extreme label skew and can outperform heavier geometry-based baselines at far lower communication cost. An ablation on the proposal variance should further confirm whether the method is robust across a range of $\sigma$ values. Future work should study adaptive switching between prototype-only and covariance-aware augmentation as the federation moves from severe class absence to milder distribution mismatch.

\bibliographystyle{splncs04}
\bibliography{references}

\end{document}
